\section{The actual common archimedean Fock quotient}
\label{sec:sr-real-matching-construction}

The archimedean comparison can be obtained from an applicable
published theta-correspondence theorem. We give the parameters,
including the character twist which makes the two targets
identical. Equality of a single lowest compact type alone
would not suffice for this purpose.

Fix an infinite place \(w\) of \(F\), with the chosen complex
embedding of \(E\). For this subsection write
\[
 \begin{gathered}
 P=a-p,\quad Q=p,\quad P'=q,\quad Q'=a'-q,\qquad
 R=a+a'-(p+q),\quad S=p+q,\\
 \kappa=m_\alpha(w)-m_\beta(w),\qquad c=\kappa/2.
 \end{gathered}
\]
Then
\begin{equation}\label{eq:sr-real-complementary-signatures}
              R=P+Q',\qquad S=Q+P',\qquad R+S=d=a+a'.
\end{equation}
The packet has \(d\) even and \(\kappa>d\).
Also \(\kappa\) is even. Indeed evaluate the restrictions of
\(\alpha,\beta\) on the \(F\)-idele whose component at \(w\)
is \(-1\) and whose other components are one.
Their infinite-type formulas give \((-1)^{m_\alpha(w)}\)
and \((-1)^{m_\beta(w)}\), whereas their quadratic restrictions
give the same sign because \(a\equiv a'\pmod2\).
Division gives \((-1)^\kappa=1\).
In particular \(c\) is an integer.
The stricter packet inequalities imply \(P,Q,P',Q'\ge0\);
they do not require \(P,Q,P',Q'\) all to be positive.

\subsection{The polynomial model and its two scalar shifts}
Choose Hermitian frames carrying the given forms and their
Cartan involutions to signatures \((P,Q)\), \((P',Q')\)
and \((R,S)\). The last frame is chosen once, for the same
skew-Hermitian space \(W\), and is used on both sides.
These frames exist by diagonalizing the positive definite
form \(x\mapsto\langle\theta x,x\rangle\) and the involution
\(\theta\). The tensor product frames carry the compatible
positive complex structures to the usual Fock structures.

Here the ordinary Fock polynomial spaces are
\[
 \mathcal P=\mathbb C[z_{ij}:1\le i\le a,\ 1\le j\le d],
 \qquad
 \mathcal P'=\mathbb C[z'_{ij}:1\le i\le a',\ 1\le j\le d].
\]
Their canonical commuting differentiated actions will be denoted
\(A_1,A_2\) and \(A'_1,A'_2\).
The normalization is the one in which the compact diagonal
vacuum weights are
\[
 \begin{array}{c|cc}
   &\text{positive block}&\text{negative block}\\ \hline
 A_1 &(R-S)/2&-(R-S)/2\\
 A_2 &(P-Q)/2&-(P-Q)/2 .
 \end{array}
\]
The primed formulas replace \(P,Q\) by \(P',Q'\).
These constants can be checked directly on a single oscillator:
differentiation of a rotation is \(z\partial_z+1/2\),
and changing its complex orientation changes its sign.
Summing over the positive and negative partner coordinates
gives exactly the displayed half-signature constants.
Off-diagonal operators follow from the corresponding
quadratic multiplication and second-derivative operators.

The chosen additive character has infinite component
\(\exp(-2\pi i\,\operatorname{Tr}_{F/\mathbb Q}x)\).
Apply the real symplectic Fourier change of variables
associated with the tensor frames, with this sign.
It carries the canonical Gaussian to a prescribed-model
Gaussian and intertwines all infinitesimal Heisenberg operators.
There is no new Gaussian-existence assumption in this step:
in diagonal real coordinates its equations are
\[
               (\partial_{x_j}+2\pi x_j)f=0,
\]
whose Schwartz solutions are precisely
\(C\exp(-\pi\sum_jx_j^2)\).
The equations with the opposite complex orientation have
only the zero Schwartz solution, because their nonzero
solutions grow like \(\exp(\pi\sum_jx_j^2)\).
Repeated creation operators give every polynomial times the
Gaussian, while multiplication and differentiation preserve
that space. Thus the smallest Heisenberg-derivative-stable
space containing all the prescribed Gaussians is precisely
the transported polynomial space. This is the source's
\(\mathrm{Pol}\), with its given Cartan metric.

The specified splitting characters add the following scalar
trace terms to the two canonical actions:
\begin{equation}\label{eq:sr-real-source-fock-shifts}
 \begin{array}{c|cc}
    &\text{first member}&\text{second member}\\ \hline
 V   & A_1-c\,\operatorname{tr}
        & A_2+\frac{m_\beta(w)}2\,\operatorname{tr}\\[2pt]
 V'  & A'_1+c\,\operatorname{tr}
        & A'_2+\frac{m_\alpha(w)}2\,\operatorname{tr}.
 \end{array}
\end{equation}
To check the factor \(1/2\), a unitary determinant \(e^{it}\)
is lifted under Hilbert's theorem 90 by \(e^{it/2}\);
the infinite-type character \((z/|z|)^m\) therefore contributes
\(e^{imt/2}\). Differentiation gives \(m\,\operatorname{tr}/2\).
On the first member use the quotient characters
\(\beta/\alpha\) and \(\alpha/\beta\); on the second use
\(\beta\) and \(\alpha\). The canonical half-signature
contribution was already included in \(A_i,A'_i\) and is
not added a second time.
The Siegel rule fixes the lift, and Egorov covariance fixes
the remaining quadratic operators, so this calculation
identifies the actual prescribed actions.

\subsection{Checking the oscillator sign in the prescribed coordinates}
\label{subsec:sr-real-oscillator-sign}

We verify the sign of the preceding vacuum weights directly;
using the conjugate oscillator would interchange the two
character-lift cases. Write
\[
 E_P=\operatorname{diag}(1_P,-1_Q),\qquad
 E_R=\operatorname{diag}(1_R,-1_S).
\]
The actual frames satisfy \(A^{\mathsf t}J_V\overline A=E_P\)
and \(C^{\mathsf t}(iJ_W)\overline C=E_R\).
Thus \(J_W\) has framed matrix \(-iE_R\), not \(iE_R\).
For framed tensor coordinates \(c=(c_{ij}),c'=(c'_{ij})\),
put \(\epsilon_i=1,-1\) on the two first blocks and
\(\eta_j=1,-1\) on the two second blocks.
The infinite symplectic form used in the model is
\[
 \operatorname{Tr}_{\mathbb C/\mathbb R}
      \left(-\sum_{ij}c_{ij}\epsilon_i(-i\eta_j)
                                      \overline{c'_{ij}}\right)
        =-2\operatorname{Im}
               \sum_{ij}\epsilon_i\eta_jc_{ij}\overline{c'_{ij}} .
\]
For each pair of signs use the real-linear coordinates
\[
 t_{ij}(c)=
   \begin{cases}
       c_{ij},&\epsilon_i\eta_j=1,\\
       \overline{c_{ij}},&\epsilon_i\eta_j=-1.
   \end{cases}
\]
Then the form is
\(-2\operatorname{Im}\sum t_{ij}(c)\overline{t_{ij}(c')}\).

There is an independent check on the additive-character sign.
With Darboux coordinates \((a,b)\), the prescribed
Schr\"odinger line is
\[
 \rho(t(a,b))f(x)=
       \exp\!\bigl(-2\pi i(t\,b\cdot x+t^2 a\cdot b/2)\bigr)
                                                          f(x+ta).
\]
Its derivative is
\(D(a,b)=\partial_a-2\pi i\,b\cdot x\);
its commutator is \(-2\pi i\) times the symplectic form.
Therefore in the \(t\)-coordinates the commutator is
\[
        [D(c),D(c')]=
            4\pi i\,\operatorname{Im}
                \sum t_{ij}(c)\overline{t_{ij}(c')}.
\]
The polynomial Heisenberg operators realizing exactly this
commutator and the decaying Gaussian are
\[
 D(c)=\sqrt{2\pi}\sum_{ij}
        \bigl(t_{ij}(c)z_{ij}
                         -\overline{t_{ij}(c)}\,\partial_{ij}\bigr).
\]
Indeed \([\partial_{ij},z_{ij}]=1\), so expansion gives
the displayed commutator. Written as \(D_t(t)\) in the
new positive complex coordinates, it also satisfies
\(D_t(it)1=iD_t(t)1\). This fixes the choice
of oscillator from the actual negative-exponential character.

Set \(s_{ij}=\epsilon_i\eta_j\).
The canonical, symmetrically ordered rotation operator
on that coordinate is
\[
              s_{ij}\frac{z_{ij}\partial_{ij}
                                  +\partial_{ij}z_{ij}}2
                 =s_{ij}(z_{ij}\partial_{ij}+1/2).
\]
For a first-member rotation in the \(i\)-th coordinate,
\(c_{ij}\) changes to \(e^{it}c_{ij}\).
Its \(t_{ij}\)-coordinate rotates by \(e^{it}\) if
\(s_{ij}=1\) and by \(e^{-it}\) if \(s_{ij}=-1\).
The preceding ordered operator has exactly that infinitesimal
Egorov action on \(D(c)\).
The \(1/2\) is the metaplectic normal-ordering constant;
it is not an extra splitting-character term.
Consequently, for a monomial \(z^\alpha\),
\[
 A_1(E_{ii})z^\alpha
   =\left(\sum_j s_{ij}(\alpha_{ij}+1/2)\right)z^\alpha,\qquad
 A_2(E_{jj})z^\alpha
   =\left(\sum_i s_{ij}(\alpha_{ij}+1/2)\right)z^\alpha.
\]
At \(\alpha=0\), these give the first-member weights
\((R-S)/2,-(R-S)/2\) and the second-member weights
\((P-Q)/2,-(P-Q)/2\), in the stated positive/negative
block order. Thus the positive-character Paul--Trapa
case really is the unprimed one.

The exact frozen-source checks are as follows.
\srcid{EvIHalfTrace.phase_line} and
\(\mathrm{dop\_apply}\) differentiate the displayed
negative-exponential line.
\srcid{RLFock.herm_frame},
\(\mathrm{exists\_heis}\), \(\mathrm{tw1}\), and
\(\mathrm{heisOp\_eq}\) give precisely the signed coordinate
change and polynomial Heisenberg operators above.
\(\mathrm{RLLieH.fock\_eq}\) identifies the symmetrized
quadratic operator with \(A_2\), and
\(\mathrm{sw\_fock}\) transposes it to \(A_1\).
Finally
\srcid{Formalism.LocalRingCardinality.DataTheta.cartan1}
and \(\mathrm{cartan2}\) are the monomial formulas just
calculated. These statements can also be checked by the
four choices of \((\epsilon_i,\eta_j)\):
the raising/lowering pairs are \((z,\partial)\)
in equal-sign blocks and \((-\partial,z)\)
in opposite-sign blocks, and their ordered diagonal
operators are \(z\partial+1/2\) and
\(-\partial z+1/2=-(z\partial+1/2)\), respectively.

\subsection{An applicable character-lift theorem}
We use Paul--Trapa, \emph{One-dimensional representations of
\(U(p,q)\) and the Howe correspondence},
J.~Funct.~Anal.~195 (2002), 129--166,
Theorem~4.1(1),(2), with the cohomological-induction normalization
of their Sections~2.3 and~2.8.
\footnote{\url{https://doi.org/10.1006/jfan.2002.3974};
author version:
\url{https://www.math.utah.edu/~ptrapa/char6_4.pdf}.}
In the positive-character case their target is
\(A_{\mathfrak q}(\det^{k/2}\otimes1)\);
in the negative-character case it is
\(A_{\mathfrak q}(1\otimes\det^{k/2})\).
The two ordered Levi signatures are respectively
\((p,q),(r-p,s-q)\) and
\((r-q,s-p),(q,p)\).
The character must occur in joint harmonics.
Their Corollary~4.3 gives nonvanishing from this occurrence.
Their square-root covers are retained; half-integral
weights must not silently be regarded as characters of
a linear unitary group.

Here is a direct check of the occurrence condition in our parameters.
Put
\[
       u=c-(R-S)/2>0,\qquad v=c+(R-S)/2>0.
\]
Both are positive integers, since \(\kappa>d\ge|R-S|\)
and \(d,\kappa\) are even.
The first compact character \(\det^c\), after subtracting
the oscillator vacuum weights, has rectangular positive
weights \(u\) on its \(P\) positive coordinates and \(v\)
on its \(Q\) negative coordinates.
These rectangles fit in the \(R\) and \(S\) partner
coordinates, because \(P\le R,Q\le S\).
Concretely products of powers of the corresponding
\(P\)-by-\(P\) and \(Q\)-by-\(Q\) minors give a nonzero
joint-harmonic highest vector.
It uses the two diagonal Fock blocks.
Every joint contraction differentiates one of the absent
off-diagonal blocks, and therefore annihilates that product.
A determinant of size zero is one, so this also covers
\(P=0\) or \(Q=0\).
For \(\det^{-c}\) on the primed side the two negative
rectangles instead occupy the off-diagonal Fock blocks
of sizes \(Q'\le R\) and \(P'\le S\).
The same contraction check proves occurrence.
Their parameters are \(k=\kappa\) and \(k=-\kappa\),
respectively; the genuine-cover congruence
\(k\equiv R-S\pmod2\) holds on both sides.

Let \(\Theta_+\) be the irreducible theta quotient of
\(\mathcal P\) for the first character \(\det^c\),
and let \(\Theta_-\) be the corresponding quotient of
\(\mathcal P'\) for \(\det^{-c}\).
The theorem just cited gives nonzero surjective quotient maps
\[
 \mathcal P\twoheadrightarrow\Theta_+,\qquad
                     \mathcal P'\twoheadrightarrow\Theta_-,
\]
on which \(A_1(X)\) and \(A'_1(X)\) act respectively by
\(c\operatorname{tr}X\) and \(-c\operatorname{tr}X\).
Thus these are quotients of the full polynomial spaces,
rather than just maps on their harmonic highest vectors.

\subsection{The two targets are the same module}
Let \(\mathfrak q\) be the theta-stable parabolic with ordered
Levi signatures
\[
                 (P,Q),\ (Q',P').
\]
For the positive side, \eqref{eq:sr-real-complementary-signatures}
makes this exactly the parabolic in the positive-character
formula, and the inducing character is
\(\det_{P,Q}^c\otimes1\).
For the negative side the first ordered block in its formula
is
\[
          (R-Q',S-P')=(P,Q),
\]
and the second is \((Q',P')\). Its parabolic is therefore
the very same \(\mathfrak q\), with inducing character
\(1\otimes\det_{Q',P'}^{-c}\).
Consequently
\[
 \Theta_+=A_{\mathfrak q}(\det_{P,Q}^c\otimes1),
 \qquad
 \Theta_-=A_{\mathfrak q}(1\otimes\det_{Q',P'}^{-c}).
\]
Cohomological induction commutes with twisting by a character
of the whole group. This follows already at the inducing
cochain complex: tensor every term by that one-dimensional
character; its restriction on the Levi is the product of
its two determinant characters, and all differentials
commute with this tensoring.
Hence
\begin{equation}\label{eq:sr-real-common-derived-module}
                   \Theta_-\otimes\det_{R,S}^c
                                     \simeq\Theta_+ .
\end{equation}
No assertion that an infinitesimal character or a compact
type uniquely determines an arbitrary module is used.

The cover identifications also agree. The second-member
metaplectic covers for the two dual pairs are controlled
by \(P-Q\equiv a\) and \(P'-Q'\equiv a'\pmod2\).
These parities coincide. Identify the covers by lifting
paths from the identity, so their differentiated actions
are identified without an extra trace shift.
Their central order-two elements act by the same sign.
The twist in \eqref{eq:sr-real-common-derived-module} is
the honest integral character \(\det^c\), so does not
change that sign. The normalized cohomological-induction
functor and its square-root Levi cover are the same for
both appearances of \(\mathfrak q\); there is consequently
no uncancelled \(\rho(\mathfrak u)\) correction.

Put \(N_w=\Theta_+\) and use
\eqref{eq:sr-real-common-derived-module} to give the primed
quotient the same codomain. We obtain onto maps \(p_w,p'_w\)
satisfying, for all complexified first and second Lie algebra
matrices \(X,Y\),
\begin{align}
 p_w(A_1(X)f)&=c\,\operatorname{tr}(X)\,p_w(f),&
 p'_w(A'_1(X)f')&=-c\,\operatorname{tr}(X)\,p'_w(f'),\label{eq:sr-real-first-lifts}\\
 N_w(Y)p_w(f)&=p_w(A_2(Y)f),&
 N_w(Y)p'_w(f')&=p'_w(A'_2(Y)f')+
                        c\,\operatorname{tr}(Y)\,p'_w(f').
 \label{eq:sr-real-second-lifts}
\end{align}
These are precisely the two lift parameters \((c,0)\)
and \((-c,c)\).
Combining \eqref{eq:sr-real-source-fock-shifts} with
\eqref{eq:sr-real-first-lifts} kills the actual first actions.
On the second actions, the primed scalar relative to \(N_w\) is
\[
                m_\alpha(w)/2-c=m_\beta(w)/2.
\]
Thus the same actual second action on both quotient maps is
\(N_w(Y)+m_\beta(w)\operatorname{tr}(Y)/2\).

As a useful independent numerical check, their joint-harmonic
compact types also agree after this twist. Writing
\(b=(P-Q)/2\) and \(b'=(P'-Q')/2\), the positive side has weights
\[
 ( (u+b)^P,b^{Q'}\ ;\ (v-b)^Q,(-b)^{P'} ).
\]
The negative side, before twisting, has
\[
 ( (b')^P,(-u+b')^{Q'}\ ;\
                    (-b')^Q,(-v-b')^{P'} ).
\]
Here \(x^j\) means that the weight \(x\) is repeated \(j\)
times. The identities \(u+b=c+b'\) and
\(v-b=c-b'\) show equality after adding \(c\) to every
primed weight. This check detects both a wrong Levi order
and a missing determinant twist.

\subsection{Compact actions, scalar endomorphisms and tensoring}
The quotient maps respect the actual compact actions as well.
Each polynomial is compact finite, so lies in a finite-dimensional
space stable under the relevant compact Lie algebra.
A Lie-algebra intertwiner on that space intertwines the
exponentials. The actual maximal compact groups are products
of connected unitary groups, and every element is an
exponential after unitary diagonalization. The first-member
action is therefore the identity on either quotient.
The second-member compact actions coincide because their
Lie actions coincide. The genuine-cover signs and the
original splitting characters have already been matched,
so this passage introduces no monodromy condition.

Each \(N_w\) is nonzero and irreducible.
It has countable complex dimension because it is a quotient
of a polynomial space in finitely many variables.
Its equivariant endomorphisms are scalar.
Here is an algebraic proof which also explains tensoring.
If a commuting endomorphism \(T\) is nonscalar, irreducibility
makes every \(T-\lambda\) invertible. For a nonzero vector \(x\),
the vectors \((T-\lambda)^{-1}x\), as \(\lambda\) ranges over
\(\mathbb C\), are linearly independent.
A finite relation, after multiplication by
\(\prod_\lambda(T-\lambda)\), would give a nonzero polynomial
in \(T\) killing \(x\); factor that polynomial into linear
factors to contradict the invertibility of each factor.
An uncountable independent set contradicts countable dimension.
Thus \(T\) is scalar.

We recall the finite-vector density argument to avoid a hidden
irreducibility assumption when tensoring different infinite places.
Let an algebra act irreducibly on a complex vector space \(M\)
with commuting endomorphisms exactly \(\mathbb C\).
For independent \(x_1,\ldots,x_t\), the set of tuples
\((A x_1,\ldots,A x_t)\) is all of \(M^t\).
Induct on \(t\). Projection to the first \(t-1\) coordinates
is onto. Its kernel in the last coordinate is either \(0\)
or \(M\). In the second case the assertion follows; in the
first it is the graph of a module homomorphism
\(M^{t-1}\to M\), hence of a scalar linear combination
of the coordinates. Applying the identity element would
make \(x_t\) dependent on the preceding vectors, a contradiction.
This proves the assertion.
Given a nonzero finite tensor sum, apply this argument to
independent vectors in its first factor to isolate a nonzero
pure tensor, and repeat in the other factors.
Irreducibility of each factor then generates every pure tensor.
Therefore the finite external tensor product is irreducible.

Take
\[
       N_\infty=\bigotimes_{w\mid\infty}N_w,\qquad
       p_\infty=\bigotimes_{w\mid\infty}p_w,\qquad
       p'_\infty=\bigotimes_{w\mid\infty}p'_w .
\]
Use the simultaneous frame transports above to view their
domains as the two actual polynomial Gaussian spaces.
Both maps are onto, since a pure tensor in the target has
a preimage obtained by choosing one preimage in each factor.
They annihilate all actual first-member Lie directions
and identify the first-member compact actions with the identity.
They intertwine the same actual second-member Lie and compact
actions. The common image is nonzero and irreducible.
These statements hold independently in every infinite-place
factor and thus also for Lie directions supported at a
single infinite place, which are the directions required
in the source's irreducibility clause.

Finally the source regards the maps as linear maps on the
whole infinite Schwartz space while imposing these properties
on its polynomial Gaussian subspace. Extend each map from that
subspace by choosing an algebraic linear complement and putting
it equal to zero on the complement.
Every asserted action and derivative remains in the polynomial
Gaussian subspace, so this extension changes none of the
properties just proved. No continuity of an arbitrary algebraic
extension is claimed or required. We have therefore produced
the actual archimedean matching datum, including its full
specified domain, rather than an abstract isomorphism of two
local theta lifts.

\subsection{Scope}
This proves the archimedean common-quotient clause for every
actual triple, using the published character-lift theorem with
its exact ranks, signatures, parity, cover and twist.
Combined with the inert, split and ramified finite constructions,
it supplies the local comparison maps simultaneously for a
fixed global triple.
The global nonzero theta-integral comparison constant, its
normalization, the two codimension-two recursion identities,
and the covolume/metric estimate remain distinct assertions.
None follows solely from the existence of these local quotients.
